% !PRINTED 586
\weberpage{586}

%<<CONT
$\equiv 1$ (mod 4), and hence $a^2+b^2\equiv 1$ (mod 4). Theorem 3.,
which moreover includes the second, is proved as follows.

Suppose that $a^2+b^2$, but not $b$, is divisible by $q$;
we then determine $b'$ from the congruence $bb'\equiv 1$ (mod $q$) and
obtain from $a^2+b^2\equiv 0$ (mod $q$)
\[
(ab')^2\equiv -1\ (\mathrm{mod}\ q).
\]
But this is impossible, since by \S.\,138, 4. the number $-1$ is a
quadratic non-residue for every prime $q$.

The field $R(i)$ is the totality of all numbers of the
form $x+yi$, where $x$, $y$ are integral or fractional rational
numbers.

\gesperrt{The numbers} $a+bi$, \gesperrt{in which} $a$ \gesperrt{and} $b$ \gesperrt{are integers
in} $R$, \gesperrt{are called the integers of the field} $R(i)$.

The product of two conjugate numbers
\[
\xi=x+yi,\quad \xi'=x-yi,
\]
that is
\[
\xi\xi'=x^2+y^2
\]
is called the norm of $\xi$ and is denoted by $N(\xi)$. The norm
of a non-vanishing number $\xi$ is a positive rational
number, and the norm of an integer is an integer. The
norm of a product or of a quotient is equal to the product
or the quotient of the norms.

\gesperrt{An integer whose norm is equal to 1 is called
a unit.}

Since $a^2+b^2$ (for integers $a$, $b$) can be $=1$ only
when $a=\pm 1$, $b=0$ or $a=0$, $b=\pm 1$, there are in
$R(i)$ only the four units
\[
+1,\quad -1,\quad +i,\quad -i.
\]
The reciprocal value of a unit is also a unit.

Two numbers, of which one arises from the other by
multiplication with a unit, are called \gesperrt{associated
numbers}.

Every complex number belongs to a system of four associated
numbers
\[
a+bi,\quad -a-bi,\quad -b+ai,\quad b-ai.
\]
The sum, difference and product of two integers are
again integers.

\gesperrt{An integer} $\alpha$ \gesperrt{is called divisible by an integer} $\beta$
\gesperrt{if a third integer} $\gamma$ \gesperrt{exists, such
that} $\alpha=\beta\gamma$.
